\section*{الفصل \ref{ch:b3:cancer-biology} --- بيولوجيا السرطان}

\begin{solution}{exo:b3:cancer-biology:1}
\omterm{def:b3:cancer-biology:genes}{المورّثة الورمية} صورة كسبية الوظيفة من مورّثة تعزز
التكاثر أو البقاء؛ ويكفي أليل طافر واحد (فهي سائدة):
مثل \emph{RAS} (بطفرة نقطية)، و\emph{MYC} (بانتقال أو تضخيم)،
و\emph{HER2}، و\emph{BCR--ABL}. أما \omterm{def:b3:cancer-biology:genes}{المورّثة الكابحة للورم} فتكبح
التكاثر أو تفرض التوقف أو الموت أو الإصلاح؛ وعلى أليليها معًا أن
يُفقدا (فهي متنحية على مستوى الخلية): مثل \emph{RB}، و\emph{TP53}، و\emph{APC}،
و\emph{BRCA1}.
\end{solution}

\begin{solution}{exo:b3:cancer-biology:2}
التكاثر المستديم: \emph{RAS}، و\emph{MYC}. وعدم الحساسية
لكوابح النمو: \emph{RB}، و\emph{CDKN2A} (p16). ومقاومة الموت:
\emph{BCL2}، و\emph{TP53}. والخلود التضاعفي: مروّج
\omterm{def:b3:cancer-biology:telomerase}{التيلوميراز}. \omterm{prop:b3:cancer-biology:oxygen}{وتولد الأوعية}: \emph{VEGF} محرَّضًا عبر HIF (بفقد
\emph{VHL}). \omterm{def:b3:cancer-biology:metastasis}{والغزو} \omterm{def:b3:cancer-biology:hallmarks}{والنقيلة}: فقد E-cadherin، و
العوامل الظهارية المتوسطية. ولااستقرار الجينوم: مورّثات \omterm{prop:b3:dna-repair:fidelity}{إصلاح عدم التوافق}،
و\emph{\omterm{def:b3:dna-repair:dsb}{BRCA}}. والاستقلاب المنفلت: \emph{MYC}، \omterm{prop:b3:cancer-biology:pathways}{ومسلك PI3K}.
\omterm{prop:b3:cancer-biology:immune}{والإفلات المناعي}: \emph{PD-L1}، وفقد MHC. والالتهاب المعزز \omterm{def:b3:cancer-biology:hallmarks}{للورم}:
إشارات \emph{NF-$\kappa$B}.
\end{solution}

\begin{solution}{exo:b3:cancer-biology:3}
الأطفال ذوو التاريخ العائلي أصيبوا بعدة أورام في العينين،
باكرًا؛ والحالات المتفرقة أصيبت \omterm{def:b3:cancer-biology:hallmarks}{بورم} واحد، متأخرًا. فاستنتج أن
حدثين لازمان في خلية: فالمرضى العائليون ورثوا
واحدًا ولم يحتاجوا إلا إلى حدث جسدي ثانٍ (متواتر بما يكفي ليقع
في عدة خلايا شبكية)، بينما احتاج المرضى المتفرقون إلى وقوع الاثنين
جسديًا في خلية واحدة --- وهو نادر ومتأخر ومفرد. ومن ثم مورّثة يجب أن تُفقد
نسختاها معًا: وهي أول \omterm{def:b3:cancer-biology:genes}{مورّثة كابحة للورم}.
\end{solution}

\begin{solution}{exo:b3:cancer-biology:4}
طافر \emph{Salmonella} المحتاج إلى الهستيدين يُزرع بلا
هستيدين مع المركب المختبَر؛ ولا تنشأ المستعمرات إلا من خلايا
عادت فيها طفرة جديدة بالعيب إلى أصله، فيقيس عددها
التطفير. ويُضاف مستخلص الكبد لأن كثيرًا من المسرطنات
غير ضارة حتى تحوّلها إنزيمات الكبد إلى مستقلَبات فعالة،
كما يحدث في الجسم؛ والبكتيريا تخلو من تلك الإنزيمات.
\end{solution}

\begin{solution}{exo:b3:cancer-biology:5}
يرتفع الحدوث $100$ ضعف بينما يتضاعف العمر: أي $2^{k-1} = 100$، ومنه $k - 1 =
\log_{2}100 = 6.6$، أي نحو سبعة أو ثمانية أحداث محددة للسرعة على
القراءة الساذجة --- وأقل إذا سُمح بالتوسع النسيلي بين الإصابات.
\end{solution}

\begin{solution}{exo:b3:cancer-biology:6}
$L = \sqrt{2Dc_{0}/q}$: فمن أجل $q = 0.01$، $\sqrt{2\times 2\times 10^{-9}
\times 0.05/0.01} = \qty{141}{\micro m}$؛ ومن أجل $q = 0.03$،
\qty{82}{\micro m}. \omterm{def:b3:cancer-biology:hallmarks}{والورم} السريع النمو يتنفس أسرع، فيستنفد
الأكسجين خلال مسافة أقصر، ويموت لبه عند حجم أصغر.
\end{solution}

\begin{solution}{exo:b3:cancer-biology:7}
$10^{11}\times 10^{-3n} < 1$ تقتضي $n > 3.7$: أي أربعة أشواط. وبعد شوطين
تبقى $10^{5}$ خلية --- أي عُشر ميليمتر مكعب، وهو أقل بكثير
من $10^{9}$ القابلة للكشف --- فيكون \omterm{def:b3:cancer-biology:hallmarks}{الورم}
``زال'' بكل اختبار بينما تبقى مئة ألف خلية؛ فيجب أن يستمر
العلاج للعدد المخطط من الأشواط، لا حتى يختفي \omterm{def:b3:cancer-biology:hallmarks}{الورم}.
\end{solution}

\begin{solution}{exo:b3:cancer-biology:8}
كلتا الآفتين تنشّط المسلك نفسه (مستقبل $\to$ \omterm{def:b3:cancer-biology:genes}{Ras} $\to$ ERK)؛
فمتى وُجدت إحداهما لم تمنح الأخرى مزيدًا من الأفضلية ولم يُنتقَ لها. و
مثبط EGFR يحجب المستقبل؛ والخلية التي تكتسب طفرة في
KRAS بعده تعيد الإشارة والمستقبل ما يزال
محجوبًا، فيُنتقى لها أثناء العلاج --- أي مقاومة بالالتفاف.
\end{solution}

\begin{solution}{exo:b3:cancer-biology:9}
يحرر E7 عامل E2F من Rb، فيدفع الخلية الظهارية المصابة عبر
\omterm{def:b3:cell-cycle-apoptosis:checkpoints}{نقطة التقييد} بلا عوامل نمو؛ ويهدم E6 بروتين \omterm{def:b3:cancer-biology:genes}{p53}، فلا يطلق
التكاثر غير الملائم ولا الأذية التي يسببها توقفًا
ولا \omterm{def:b3:cell-cycle-apoptosis:apoptosis}{استماتة}. فتُوفَّر سمتان دفعةً واحدة، لكن
الباقيات --- الخلود، \omterm{def:b3:cancer-biology:metastasis}{والغزو}، واللااستقرار --- تحتاج إلى طفرات
جسدية أخرى، تستغرق عقودًا لتتراكم في الخلايا المصابة
إصابة مستمرة. واللقاح يثير أجسامًا مضادة معادلة تجاه
القفيصة ويمنع الخمج؛ فإذا أُعطي قبل التعرض لم توجد
خلية مصابة قط لتتحول.
\end{solution}

\begin{solution}{exo:b3:cancer-biology:10}
بعد الإصابة $i$ تنمو النسيلة بعامل $g$، فيكون $g^{i}$ ضعف
الخلايا معرَّضًا للإصابة التالية؛ واحتمال المتتالية الكاملة
في السلالة المنحدرة من خلية أصلية واحدة هو $g\cdot g^{2}\cdots$
--- ومن أجل $g$ ثابتة لكل إصابة يصير الخطر التراكمي $N(ut)^{k}
g^{k-1}$ حتى عوامل تتوقف على الترتيب، ويحتفظ الحدوث
بالأسّ $t^{k-1}$ بمعامل أكبر قدره $g^{k-1}$. وإذا نمت التوسعات
نفسها مع الزمن (أي نسائل نامية أسّيًا)، ازداد المنحني حدةً،
فيُنتَج ميل قدره $6$ بأقل من سبع مورّثات
قائدة: فالمقدار $k$ في أرميتاج--دول حد أعلى لعدد
القائدات المحددة للسرعة، والأورام المسلسلة تُظهر اثنتين إلى ثماني.
\end{solution}

\begin{solution}{exo:b3:cancer-biology:11}
(أ) الدواء A وحده يواجه $\mu_{A}N = 100$ خلية مقاومة موجودة سلفًا:
ومنه $P_{0} = e^{-100} \approx 0$؛ فتنجو النسيلة المقاومة له و
تعيد النمو خلال الأسابيع الثمانية عشر من A، ومتى بدأ B كانت
جمهرة كبيرة من جديد، فيواجه B أيضًا $\mu_{B}N' \gg 1$ --- أي يخفق
التتابع مرتين. (ب) ومعًا: تنشأ خلية مقاومة للاثنين بمعدل
$10^{-14}$ لكل انقسام، و$\mu_{A}\mu_{B}N = 10^{-5}$، و$P_{0} =
0.99999$. فالتوليفة من البداية، حين يكون $N$ أصغر ما يكون، هي
القاعدة لأن جداء معدلين صغيرين هو ما يجعل المقاومة
غير مرجحة.
\end{solution}

\begin{solution}{exo:b3:cancer-biology:12}
مع خلايا أكثر بمقدار $100$ ضعف وبالقيمتين $u$ وَ$k$ نفسيهما، كان الخطر التراكمي
$N(ut)^{k}$ ليكون أعلى $100$ ضعف --- فينبغي أن يموت كل فيل من
\omterm{def:b3:cancer-biology:hallmarks}{السرطان} صغيرًا. وهو لا يموت، فلا بد أن تختلف $u$ أو $k$: فعشرون نسخة من
\emph{TP53} تجعل الاستجابة الاستماتية للأذية أقوى وتزيل
الخلايا الطافرة قبل أن تتوسع (فتخفض $u$ الفعالة لتلك
الخطوة)؛ كما أن للحيوانات الكبيرة معدلات طفرة أدنى في الخلية (بإصلاح
أفضل، وبانقسامات أقل في الخلية كل سنة لأن خلاياها تدور
أبطأ)، وقد تحتاج إلى إصابات أكثر (بطبقة كابحة إضافية).
فقد تطور كبح \omterm{def:b3:cancer-biology:hallmarks}{السرطان} مع حجم الجسم.
\end{solution}

\begin{solution}{pb:b3:cancer-biology:1}
\textbf{1.} $10^{12}/10^{3} = 10^{9}$ خلية؛ و$\log_{2}10^{9} \approx 30$
مضاعفة.
\textbf{2.} $30\times 100 = 3000$ يوم، أي نحو $8$ سنوات.
\textbf{3.} \qty{1}{kg} $\approx 10^{12}$ خلية: أي عشر مضاعفات أخرى،
و$1000$ يوم، أي دون ثلاث سنوات --- أي ثلث الزمن حتى الكشف.
\textbf{4.} لا بد أن زمن المضاعفة الباكر كان نحو $60$ يومًا، أو
أن \omterm{def:b3:cancer-biology:hallmarks}{الورم} نما أسرع وهو صغير وبطؤ مع نموه (وهو النمط المعتاد،
إذ يخفق الإمداد والحيّز).
\textbf{5.} دورة من $2$ يوم كانت لتضاعف الجمهرة كل $2$
يوم؛ وزمن مضاعفة قدره $100$ يعني أن النمو الصافي $2\,\%$ من
معدل الولادة: أي إن كل خلية تولد يقابلها تقريبًا خلية تموت، أو إن القليل
منها في الدورة --- فالولادة والموت متوازنان تقريبًا، والنمو هو الفرق
الصغير.
\textbf{6.} الأدوية تقتل الخلايا في الدورة، \omterm{def:b3:cancer-biology:hallmarks}{فالورم} الذي نسبة قليلة
منه في الدورة يفقد نسبة أقل في الشوط؛ لكن \omterm{def:b3:cancer-biology:hallmarks}{الورم} السريع، وقد
فقد أكثر، يعيد النمو أيضًا بين الأشواط. وفي المحصلة تكون الأورام
السريعة (الابيضاضات، واللمفومات، \omterm{def:b3:cancer-biology:hallmarks}{وسرطان} الخصية) هي التي تشفيها
المعالجة الكيميائية.
\textbf{7.} النسبة $300$ على نسبة عمر $80/30 = 2.67$: ومنه $k - 1 =
\ln 300/\ln 2.67 = 5.8$، و$k \approx 7$.
\textbf{8.} $N(ut)^{k} = 10^{8}\times (7\times 10^{-5})^{7} \approx
10^{-21}$: وهو عبث في مواجهة خطر مدى الحياة قرب $1$ من $3$. والناقص:
التوسع النسيلي الذي يضاعف الخلايا المعرَّضة بعد كل إصابة، ومعدلات
الطفرة المرفوعة باللااستقرار، وانقسامات أكثر بكثير (فالخلايا الجذعية تدور
عقودًا)، وقائدات حقيقية أقل.
\textbf{9.} الضرب في $g^{k-1} = 100^{6} = 10^{12}$ يعطي $10^{-9}$
--- وهو ما يزال ضئيلًا بسبعة أحداث؛ ومع أربعة أحداث وتوسعات،
$10^{8}\times (7\times 10^{-5})^{4}\times 10^{6} \approx 2\times
10^{-3}$، وهي الرتبة الصحيحة. فالصورة الحقيقية قائدات قليلة وتوسعات
كبيرة.
\textbf{10.} $(ut)^{k-1}/(ut)^{k} = 1/(ut) = 1/(10^{-6}\times 40) =
2.5\times 10^{4}$: أي إن حدوث الحاملة عند الأربعين عشرات آلاف
الأضعاف من الحدوث المتفرق --- فالسرطانات الوراثية تضرب صغارًا.
\textbf{11.} حدث واحد أسرع عشر مرات: الحدوث $\times 10$؛ وحدثان:
$\times 100$. والخطر المرصود البالغ خمسة عشر إلى خمسة وعشرين ضعفًا يوحي
بأن الدخان يعمل على أكثر من خطوة.
\textbf{12.} عند الإقلاع تعود $u$ للأحداث التي لم تقع بعد،
فيكف الحدوث عن الصعود بتلك الحدة؛ لكن الإصابات
المتراكمة سلفًا في الخلايا تبقى، فتكون خلايا المقلع أقرب إلى
\omterm{def:b3:cancer-biology:hallmarks}{السرطان} من خلايا من لم يدخن قط ويدوم الخطر الزائد، مجمَّدًا
لا ممحوًا.
\textbf{13.} $L = \sqrt{2\times 2\times 10^{-9}\times 0.05/0.03} =
\qty{82}{\micro m}$.
\textbf{14.} مؤكسجة تمامًا حتى قطر $2L \approx
\qty{160}{\micro m}$. وللكروية البالغة \qty{2}{mm} حافة حية قدرها
\qty{82}{\micro m}: فالنسبة الحية $1 - (918/1000)^{3} = 0.23$.
\textbf{15.} خلايا المنتصف تبعد \qty{75}{\micro m} عن وعاء، أي داخل
$L$ بقليل --- فهي مؤكسجة مبدئيًا، لكن أوعية \omterm{def:b3:cancer-biology:hallmarks}{الورم} تحمل أكسجينًا أقل
ويكون جريانها متقطعًا، فتكون خلايا المنتصف ناقصة الأكسجة.
والخلايا ناقصة الأكسجة أشد مقاومة للإشعاع بضعفين إلى ثلاثة
(فالأكسجين لازم لجعل الأذية دائمة)، ونقصُ الأكسجة يطلق
\omterm{def:b3:cell-cycle-apoptosis:apoptosis}{استماتة} معتمدة على \omterm{def:b3:cancer-biology:genes}{p53}، فينتقي الخلايا التي فقدت \omterm{def:b3:cancer-biology:genes}{p53}.
\textbf{16.} تنصيف الكثافة يباعد الأوعية بعامل $\sqrt{2}$:
فتصير \qty{212}{\micro m} بعضها عن بعض، وتبعد خلايا المنتصف \qty{106}{\micro m}، أي وراء
$L$: فتموت أو، إن نجت ناقصة الأكسجة، صارت أشد عدوانية وأكثر
توليدًا للأوعية.
\textbf{17.} $30/2 = 15$ ضعفًا من الغلوكوز. فتركّز الأورام
نظائر الغلوكوز، ويضيئها مسح
بالفلورودي أوكسي غلوكوز الباعث للبوزيترونات.
\textbf{18.} الأورام المجوَّعة تتقلص إلى الحجم المحدود بالانتشار و
تدوم، راقدةً ناقصة الأكسجة، أو تستولي على أوعية قائمة؛ ونقصُ الأكسجة
ينتقي نسائل عدوانية. لكن تشذيب أوعية \omterm{def:b3:cancer-biology:hallmarks}{الورم} الفوضوية
يخفض الضغط الخلالي ويسوّي الجريان، فتنفذ المعالجة الكيميائية
أفضل --- فتنجع التوليفة حيث لا ينجع أي منهما
وحده.
\textbf{19.} $10^{9}\times 10^{-2n} < 1$: ومنه $n > 4.5$، أي خمسة أشواط.
\textbf{20.} $21$ يومًا $= 0.21$ مضاعفة، فإعادة النمو $\times 1.16$؛ والعامل
الصافي في الدورة $0.0116$؛ و$\ln 10^{-9}/\ln 0.0116 = 4.65$: أي خمس
دورات أيضًا.
\textbf{21.} $10^{-7}\times 10^{9} = 100$ خلية مقاومة متوقعة؛
و$P_{0} = e^{-100} \approx 0$: فالمقاومة يقينية.
\textbf{22.} $\mu_{1}\mu_{2}N = 10^{-14}\times 10^{9} = 10^{-5}$؛
و$P_{0} = 0.99999$. وعند $10^{12}$ خلية، $10^{-2}$، و$P_{0} = 0.99$.
\textbf{23.} $10^{6}$ خلية: فبدواء واحد $\mu N = 0.1$، و$P_{0} = 0.90$؛ وبدواءين
$10^{-8}$، و$P_{0} \approx 1$. فعالج حين تكون الجمهرة صغيرة،
وعالج بالتوليفات.
\textbf{24.} الخلايا التائية تهاجم مستضدات جديدة كثيرة دفعةً واحدة، فلا تفلت منها
طفرة واحدة؛ ويقتضي الإفلات فقد عرض المستضدات
نفسه (MHC، و$\beta_{2}$-ميكروغلوبولين) أو كبت الخلايا التائية، وهو
صنف مختلف وأندر من الأحداث. والحمل الطفري العالي نفسه الذي
يجعل \omterm{def:b3:cancer-biology:hallmarks}{الورم} يقاوم مثبط الكيناز يقينًا يمنحه مستضدات جديدة كثيرة:
فالأورام المعوزة \omterm{prop:b3:dna-repair:fidelity}{إصلاح عدم التوافق} والمحدثة بالضوء فوق البنفسجي وبالدخان
تستجيب أحسن استجابة.
\textbf{25.} نحو $8$ سنوات حتى الكشف؛ \omterm{def:b3:cancer-biology:hallmarks}{والورم} الحي عمقه \qty{82}{\micro m}
حول وعاء؛ و$P_{0} = 0.99999$ أن توليفة الدوائين
لا تواجه خلية مقاومة في \omterm{def:b3:cancer-biology:hallmarks}{ورم} قدره \qty{1}{cm^3}.
\end{solution}
